\chapter*{Abstract}
\addcontentsline{toc}{chapter}{Abstract}

Retrieval-Augmented Generation (RAG) enables natural-language access to organisational data, but source-level authorisation can be lost as information passes through retrieval, relations, conversation state, model context, generation, and delivery. This thesis implements and evaluates a modular RAG system over a fictitious but realistic product-development workbook. The system retains field sensitivity and provenance in structured entity representations and compares two context strategies: role-based projection before generation and a diagnostic condition that deliberately keeps labelled mixed-sensitivity evidence model-visible.

Eight attack families evaluate direct and multi-turn extraction, access revocation, relation and indexed-record inference, poisoned-context integrity, triggered indirect prompt injection, and retrieval-mediated numeric disclosure. The evaluation separates internal exposure, raw generation, control action, delivered output, and authorised task success. It also distinguishes descriptive observations of the historical and hardened packages from matched and supplemental component studies.

The structurally projected path performed strongly under the evaluated direct-confidentiality conditions, whereas substantial baseline failures occurred when protected evidence remained available to the model. The final hardened package recorded no unauthorised attack-specific outcome in the finite evaluation matrix, a substantial descriptive improvement over the historical implementation. Targeted component studies further showed that individual controls can eliminate specific attack outcomes under matched conditions; for example, prompt-injection protection reduced a reproduced integrity failure from 300/300 to 0/300 cases.

The results indicate that access restrictions are easier to maintain when they are enforced structurally before generation. Deliberately model-visible protected evidence requires additional, observable controls at multiple pipeline stages, and successful delivery filtering does not remove upstream exposure. Package comparison is descriptive; matched and supplemental evidence supports attribution. Finite zero-leakage evidence does not establish general security; the architecture and stage-aware evaluation framework provide a basis for broader testing across policies, datasets, models, and attack conditions.

\clearpage
